% Part: sets-functions-relations
% Chapter: induction
% Section: induction

\documentclass[../../../include/open-logic-section]{subfiles}

\begin{document}

\olfileid{sfr}{ind}{int}

\olsection{في الاستقراء}

\begin{explain}
الاستقراء طريق مألوف من طرق البرهان في الرياضيات والمنطق. ولعلك عرفته من
الاستقراء على الأعداد الطبيعية؛ غير أن الاستقراء يجري أيضًا في أصناف
كثيرة من الأشياء الرياضية والمنطقية.

والغرض منه إثبات دعوى كلية، أي إثبات خاصية لكل شيء من نوع معين. وتظهر
منفعته خاصة حين يتعذر إدخال الأشياء كلها في البرهان مباشرة. وإنما يصح
هذا الطريق في الأشياء المنشأة على نظام معلوم: فإذا كان كل عنصر من
المجموعة أصلًا، أو مبنيًا من أصول في عدد متناه من الخطوات، جرى عليها
الاستقراء. وبيان ذلك أدق بما يأتي:
\end{explain}

\begin{defn}
لتكن $f$ دالة $n$-موضعية. ونقول إن المجموعة $S$ \emph{مغلقة} تحت $f$
إذا وفقط إذا كان $f(x_1,\dots,x_n) \in S$ كلما كانت
$x_1,\dots,x_n \in S$ وكانت $f$ معرفة عند هذه المدخلات.
\end{defn}

\begin{explain}
فمجموعة الأعداد الطبيعية ($\mathbb{N}$)، مثلًا، مغلقة تحت دالة التالي
$f(x) = x+1$؛ إذ إذا كان $x$ طبيعيًا كان $x+1$ طبيعيًا. وليست مغلقة،
في المقابل، تحت دالة السابق $f(x) = x-1$؛ لأن $0$ طبيعي، وأما $-1$
فليس بطبيعي.
\end{explain}

\begin{defn}
نقول إن الخاصية $P$ \emph{محفوظة تحت} الدالة $n$-موضعية $f$ إذا كان،
لكل مدخلات $a_1,\dots,a_n$ من مجال $f$، صدق
$P(a_1),\dots,P(a_n)$ مستلزمًا لصدق $P(f(a_1,\dots,a_n))$.
\end{defn}

\begin{explain}
خاصية «كون العدد زوجيًا» محفوظة تحت الدالة $f(x) = x+2$؛ لأن
$x$ إذا كان زوجيًا كان $x+2$ زوجيًا. وأما دالة التالي فلا تحفظ هذه
الخاصية، إذ يكون $x+1$ فرديًا متى كان $x$ زوجيًا.

وإنما يصح الاستقراء على الأعداد الطبيعية لأن كل عدد طبيعي يُنال من
$0$ بتكرار دالة التالي عددًا متناهيًا من المرات. وهذه هي السمة العامة
لكل مجموعة يصح فيها الاستقراء.
\end{explain}

\begin{defn}[الاستقراء على $\Nat$]
إذا اتصف $0$ بخاصية $P$، وكانت $P$ محفوظة تحت دالة التالي، اتصف كل
عدد طبيعي بالخاصية $P$.
\end{defn}

\begin{ex} إن مجموع الأعداد الطبيعية من $0$ إلى $n$ يساوي $\frac{n(n+1)}{2}$. \\

\textbf{حالة الأساس.} المجموع من $0$ إلى $0$ هو $0=
\frac{0\cdot 1}{2}$.\\

\textbf{خطوة الاستقراء.} افترض صدق الخاصية عند $k$، ونثبت صدقها عند
$k+1$. أي نفترض صدق $P(k)$، وهو أن
\[ 0 + 1 + \ldots + k = \frac{k(k+1)}{2} \]
ثم نثبت صدق $P(k+1)$، وهو أن
\[ 0 + 1 + \ldots + k + (k+1) = \frac{(k+1)(k+2)}{2} \]
مستعينين بفرض صدق $P(k)$.

\begin{align*}
0 + 1 + \ldots + k &= \frac{k(k+1)}{2} \tag{فرض الاستقراء} \\
0 + 1 + \ldots + k + (k+1) &= \frac{k(k+1)}{2} + (k+1)
\tag{إضافة $k+1$ إلى الطرفين} \\
&= \frac{k(k+1) + 2(k+1)}{2} \\
&= \frac{k^2 + k + 2k +2}{2} \\
&=\frac{(k+1)(k+2)}{2}
\end{align*}
فتمت خطوة الاستقراء، وثبتت الدعوى.
\end{ex}

\begin{explain}
وأما !!{formula}s فأصولها !!{formula}s الذرية، وتُنشأ !!{formula}s المركبة
بضم !!{formula}s أقل تعقيدًا بواسطة الروابط. فلكل رابط دالة على
!!{formula}s، ترد !!{formula} جديدة عند تطبيق الرابط على مدخلاته. ومن
ذلك الدالة التي تقرن اثنتين من !!{formula}s، هما $!A$ و$!B$، ويمكن كتابتها
$\mathcal E _\land (!A, !B) = !A \land !B$. ولنا أن نسمي هذه وأمثالها
\emph{دوال بناء الصيغ}. ومجموعة !!{formula}s مغلقة تحتها؛ فإذا كان كل من
$!A$ و$!B$ من !!{formula}s، كان كل من $\lnot !A$ و$!A \land !B$،
وغيرهما، صيغة كذلك.
\end{explain}

\begin{defn}[الاستقراء على !!{formula}s]
إذا اتصفت كل !!{formula} ذرية بالخاصية $P$، وكانت $P$ محفوظة تحت دوال
بناء !!{formula}، اتصفت كل !!{formula} بالخاصية $P$.
\end{defn}

\begin{explain}
يفيد هذا التعريف منهاج البراهين الاستقرائية على !!{formula}s. فإذا أردنا
أن نبين اتصاف كل !!{formula} بالخاصية $P$، رتبنا البرهان هكذا:\\

\textbf{حالة الأساس.} لتكن $!A$~!!{formula} ذرية. [\ldots]
فيثبت اتصاف $!A$ بالخاصية $P$.

\textbf{خطوة الاستقراء.} لتكن $!A$ و$!B$ من !!{formula}s، ولتكن
الخاصية $P$ صادقة عليهما.

حالة $\lnot$: [\ldots] فيثبت اتصاف $\lnot !A$ بالخاصية $P$.

حالة $\land$: [\ldots] فيثبت اتصاف $!A \land !B$ بالخاصية $P$.

حالة $\lor$: [\ldots] فيثبت اتصاف $!A \lor !B$ بالخاصية $P$.

حالة $\lif$: [\ldots] فيثبت اتصاف $!A \lif !B$ بالخاصية $P$.

حالة $\liff$: [\ldots] فيثبت اتصاف $!A \liff !B$ بالخاصية $P$.

حالة $\lforall[]$: [\ldots] فيثبت اتصاف $\lforall[x] !A$ بالخاصية $P$.

حالة $\lexists[]$: [\ldots] فيثبت اتصاف $\lexists[x] !A$ بالخاصية $P$. \\

ومن ثم تتصف كل !!{formula} بالخاصية $P$.
\end{explain}

\end{document}
